\documentclass[12pt]{article}
\usepackage{latexblog}

% ============================================================
% Blog Metadata
% ============================================================
\blogtitle{关于随机数}
\blogdate{2018-12-03}
\blogtags{刨根问底, 算法}
\bloglang{zh}
\blogsummary{}

\begin{document}
\makeblogtitle

随机数顾名思义就是你无法确定的一个数（但是你可以设定一个范围），就好比彩票摇号一样，所有可能的组合是知道的，但是到底会摇出个什么数字出来，谁都不知道。否则我早就买彩票去了😂 那随机数是怎么生成出来的？

\section{随机数的定义}\label{ux968fux673aux6570ux7684ux5b9aux4e49}

引用维基百科，

\begin{quote}
根据密码学原理，随机数的随机性检验可以分为三个标准：
\end{quote}

\begin{quote}
\begin{itemize}
\tightlist
\item
  统计学伪随机性。统计学伪随机性指的是在给定的随机比特流样本中，1的数量大致等于0的数量，同理，``10''\,``01''\,``00''\,``11''四者数量大致相等。类似的标准被称为统计学随机性。满足这类要求的数字在人类``一眼看上去''是随机的。
\item
  密码学安全伪随机性。其定义为，给定随机样本的一部分和随机算法，不能有效的演算出随机样本的剩余部分。
\item
  真随机性。其定义为随机样本不可重现。实际上衹要给定边界条件，真随机数并不存在，可是如果产生一个真随机数样本的边界条件十分复杂且难以捕捉（比如计算机当地的本底辐射\footnote{本体辐射是指人类生活环境本来存在的辐射，主要包括宇宙射线和自然界中天然放射性核素发出的射线。}波动值），可以认为用这个方法演算出来了真随机数。但实际上，这也只是非常接近真随机数的伪随机数，一般认为，无论是本地辐射、物理噪音、抛硬币\ldots\ldots 等都是可被观察了解的，任何基于经典力学产生的随机数，都只是伪随机数。
\end{itemize}
\end{quote}

\begin{quote}
相应的，随机数也分为三类：
\end{quote}

\begin{quote}
\begin{itemize}
\tightlist
\item
  伪随机数：满足第一个条件的随机数。
\item
  密码学安全的伪随机数：同时满足前两个条件的随机数。可以通过密码学安全伪随机数生成器计算得出。
\item
  真随机数：同时满足三个条件的随机数。
\end{itemize}
\end{quote}

\section{Linux系统中的随机数设备}\label{linuxux7cfbux7edfux4e2dux7684ux968fux673aux6570ux8bbeux5907}

Linux以及一些类Unix系统中有随机数的特殊文件，一般如下：

\begin{itemize}
\tightlist
\item
  /dev/random :提供基于当前系统熵池\footnote{指设备驱动程序或其它来源的背景噪声计算出来的某种结果}的真随机数
\item
  /dev/urandom:是非阻塞的随机数生成器
\end{itemize}

两者都是CSPRNG\footnote{Cryptographically Secure Pseudorandom Number Generator，加密安全的伪随机数生成器}，可以使用以下命令来输出：

\begin{Shaded}
\begin{Highlighting}[]
\FunctionTok{od} \AttributeTok{{-}An} \AttributeTok{{-}N1} \AttributeTok{{-}i}\NormalTok{ /dev/random}
\end{Highlighting}
\end{Shaded}

\section{一些伪随机数生成算法}\label{ux4e00ux4e9bux4f2aux968fux673aux6570ux751fux6210ux7b97ux6cd5}

\subsection{平方取中法}\label{ux5e73ux65b9ux53d6ux4e2dux6cd5}

这个算法比较简单，由冯·诺伊曼在1946年提出。 算法步骤如下：

\begin{itemize}
\tightlist
\item
  选择一个 \({\displaystyle m}\) 位数 \({\displaystyle N_{i}}\) 作为种子
\item
  计算 \({\displaystyle N_{i}^{2}}\)
\item
  若 \({\displaystyle N_{i}^{2}}\)不足 \({\displaystyle 2m}\)个位，在前补0。在这个数选中间 \({\displaystyle m}\)个位的数，即 \({\displaystyle 10^{\lfloor {\frac {m}{2}}\rfloor +1}} {\displaystyle 10^{\lfloor {\frac {m}{2}}\rfloor +1}}\)至 \({\displaystyle 10^{\lfloor {\frac {m}{2}}\rfloor +m}} {\displaystyle 10^{\lfloor {\frac {m}{2}}\rfloor +m}}\)的数，将结果作为 \({\displaystyle N_{i+1}}\)
\end{itemize}

\subsection{线性同余法}\label{ux7ebfux6027ux540cux4f59ux6cd5}

这个算法根据递归公式计算:

\[
X_{n+1}=\left(aX_{n}+c\right)~~{\bmod {~}}~m
\]

Java中的Random类就是使用就是这种算法。但这个不是密码学安全的随机数算法，如果要生成密码学安全的随机数，需要使用SecureRandom类来生成。

\subsection{Blum Blum Shub}\label{blum-blum-shub}

采用如下的递归公式计算：

\[
x_{n+1}=x_{n}^{2}{\bmod  M}
\]

其中：\(M=p\cdot q\)是两个大素数p和q的乘积

例如令\({\displaystyle p=11}\), \({\displaystyle q=19}\), \({\displaystyle s=3}\)，则：

\begin{enumerate}
\tightlist
\item
  \({\displaystyle x_{0}=3^{2}{\bmod 209}=9}\)
\item
  \({\displaystyle x_{1}=9^{2}{\bmod 209}=81}\)
\item
  \({\displaystyle x_{2}=81^{2}{\bmod 209}=82}\)
\item
  \({\displaystyle x_{3}=82^{2}{\bmod 209}=36}\)
\item
  \ldots{}
\end{enumerate}

除此之外，还有一些其他的随机数算法，便不过多介绍。

参考:

\begin{itemize}
\tightlist
\item
  \href{http://www.2uo.de/myths-about-urandom/}{Myths about /dev/urandom}
\item
  \href{http://www.cnblogs.com/Geometry/archive/2011/01/25/1944582.html}{一个生成伪随机数的超级算法}
\end{itemize}


\end{document}
